\documentclass[11pt]{article}

\usepackage[margin=1in]{geometry}
\usepackage[T1]{fontenc}
\usepackage{lmodern}
\usepackage{amsmath,amssymb}
\usepackage{microtype}
\usepackage{array,booktabs,longtable}
\usepackage{hyperref}

\hypersetup{hidelinks}
\setlength{\parindent}{0pt}
\setlength{\parskip}{0.45em}

\title{\vspace{-1.5em}\large Math 277 Project Proposal\\[0.5em]
  \LARGE Representation-Space Optimization\\
  for Feature Learning}
\author{Jacob Cai \qquad Jiayi Liwei}
\date{November 13, 2026}

\begin{document}
\maketitle
\vspace{-1.5em}

\section*{Research question and connection to the course}

Building on the suggested project ``Representation metrics for feature
learning,'' we will investigate how optimization geometry affects neural
network representations. Our central question is:
\emph{Does a representation-space optimization rule produce more task-effective
feature motion than ordinary parameter gradient descent?}
Using one small residual architecture and one synthetic dataset, we will
investigate whether the Jacobian spectrum explains differences in learning
and whether any improvement persists after accounting for computational cost.
The project connects natural-gradient methods and pullback geometries
\cite{martens2020,arbel2020} with Jacobians, tangent kernels, and
feature-learning dynamics.

\section*{1. Representation map}

For inputs $x\in\mathbb R^2$, we will use a two-block residual network
\[
 h_0=x,\qquad
 h_{\ell+1}=h_\ell+\tfrac12 B_\ell\tanh(A_\ell h_\ell+b_\ell),
 \quad \ell=0,1,
\]
where $A_\ell\in\mathbb R^{8\times2}$,
$B_\ell\in\mathbb R^{2\times8}$, and $b_\ell\in\mathbb R^8$ are the
trainable parameters $\theta$. Define
$R_X(\theta)=\operatorname{vec}(h_2(x_1),\ldots,h_2(x_n))=Z_\theta$
and $J_\theta=DR_X(\theta)$. The final representations are two-dimensional,
and the network has 80 trainable parameters.

\section*{2. Metric and objective}

We choose the Euclidean representation metric
$\|\delta Z\|^2=\sum_i\|\delta z_i\|_2^2$ and a Euclidean parameter
penalty. With labels $y_i\in\{-1,1\}$, both methods use the same logistic
objective
$\Phi(Z)=n^{-1}\sum_i\log(1+\exp(-y_i w^{\mathsf T}z_i))$
and fixed readout $w=(1,0)^{\mathsf T}$.

\section*{3. Update rule and mathematical analysis}

Writing $g=\nabla_Z\Phi(Z_\theta)$, we will derive the damped pullback
direction by matching the desired representation velocity $-g$:
\begin{align*}
 u_\lambda
 &=\underset{u}{\operatorname{argmin}}\;
       \left\{\tfrac12\|J_\theta u+g\|_2^2
       +\tfrac\lambda2\|u\|_2^2\right\}\\
 &=-\bigl(J_\theta^{\mathsf T}J_\theta+\lambda I\bigr)^{-1}
       J_\theta^{\mathsf T}g,\qquad \lambda>0.
\end{align*}
Thus $J_\theta^{\mathsf T}J_\theta$ is the possibly singular pullback metric;
adding $\lambda I$ makes it positive definite. We will compare
$\theta\leftarrow\theta+\eta u_\lambda$
with ordinary gradient descent. Singular value decomposition will explain
the local representation factors $\sigma_i^2/(\sigma_i^2+\lambda)$ for
the pullback update versus $\sigma_i^2$ for gradient descent. We will verify
these first-order predictions numerically and establish the local descent
property of the damped direction.

\section*{4. Planned experiment}

We will first check that small numerical updates agree with the predicted
first-order feature motion, then run the controlled comparison below.
The sample sizes are initial choices that may be refined after this pilot.

\begingroup
\small
\setlength{\LTpre}{0.25em}
\setlength{\LTpost}{0.25em}
\renewcommand{\arraystretch}{1.12}
\begin{longtable}{@{}>{\raggedright\arraybackslash}p{0.21\textwidth}
  >{\raggedright\arraybackslash}p{\dimexpr0.79\textwidth-2\tabcolsep\relax}@{}}
\toprule
\textbf{Component} & \textbf{Initial experimental plan}\\
\midrule
\endfirsthead
\toprule
\textbf{Component} & \textbf{Initial experimental plan (continued)}\\
\midrule
\endhead
\bottomrule
\endfoot
Data & Noisy two moons; initially 256 training, 256 validation, and 2,000
test examples. Each paired comparison uses identical data splits.\\
\addlinespace
Methods & Full-batch gradient descent versus the damped pullback update,
using the architecture above and the same fixed readout and objective.\\
\addlinespace
Repetitions & At least three paired random seeds, sharing initialization
within each pair. Report mean performance and variability.\\
\addlinespace
Tuning & Select learning rates separately for each method and a fixed
$\lambda>0$ per pullback run, using validation data and equal tuning budgets.
Test data are reserved for evaluation.\\
\addlinespace
Primary comparisons & Plot training loss and test accuracy against both
\textbf{iteration count} and \textbf{elapsed training time}.\\
\addlinespace
Spectral explanation & At initial, intermediate, and final checkpoints,
record Jacobian singular values and gradient alignment with the corresponding
representation directions; relate these to the observed learning dynamics.\\
\addlinespace
Feature quality & Visualize the point clouds and evaluate linear probes on
initial and learned features using the same training-based normalization
and regularization-selection protocol.\\
\addlinespace
Computing and timing & Begin with a CPU pilot. Run both methods on the same
hardware, device, and numerical precision. Include required Jacobian and
linear-solve costs; record extra diagnostic costs separately.\\
\end{longtable}
\endgroup

For $n=256$, the Jacobian is $512\times80$ and the pullback step requires
an $80\times80$ linear solve, making a CPU pilot a reasonable starting point.
GPU use is optional and will be guided by measured runtime. As a secondary
diagnostic, we will compare progress at matched cumulative feature displacement,
defined as the sum of RMS changes in final representations between successive
updates on a fixed input set.

\section*{Scope and expected outcomes}

The core deliverables are the metric-to-update derivation, local spectral
analysis, and a reproducible numerical comparison. Both improvements and
limitations, such as iteration savings offset by computational cost, will
be informative.
\textbf{Only after these components are complete}, we may add a nuisance
feature whose label correlation changes at test time to explore distribution
shift \cite{chen2023}. Shifted test data will not be used for tuning.
The initial scope is the endpoint metric and the single experimental setting
above.

\renewcommand{\refname}{References}
\begin{thebibliography}{9}
\setlength{\itemsep}{0.25em}

\bibitem{martens2020}
J.~Martens.
\newblock \href{https://www.jmlr.org/papers/volume21/17-678/17-678.pdf}
{New Insights and Perspectives on the Natural Gradient Method}.
\newblock \emph{Journal of Machine Learning Research}, 21(146):1--76, 2020.

\bibitem{arbel2020}
M.~Arbel, A.~Gretton, W.~Li, and G.~Mont\'ufar.
\newblock \href{https://openreview.net/pdf?id=Hklz71rYvS}
{Kernelized Wasserstein Natural Gradient}.
\newblock \emph{ICLR}, 2020.

\bibitem{chen2023}
Y.~Chen, W.~Huang, K.~Zhou, Y.~Bian, B.~Han, and J.~Cheng.
\newblock \href{https://papers.neurips.cc/paper_files/paper/2023/hash/d73d5645ddbb9ada6c862116435574f6-Abstract-Conference.html}
{Understanding and Improving Feature Learning for Out-of-Distribution
Generalization}.
\newblock \emph{NeurIPS}, 2023.

\end{thebibliography}

\end{document}
